\documentclass[a4paper,12px]{article}
\usepackage{graphicx}
\usepackage[english]{babel}
\usepackage{fancyhdr}
\usepackage{lastpage}
\usepackage{xifthen}
\usepackage[linesnumberedhidden, titlenotnumbered]{algorithm2e}
\usepackage{lipsum}
\usepackage{hyperref}
\usepackage{array}
\usepackage{tabularx}

\usepackage{minted}
\usepackage{caption}
\usepackage{amssymb}

\pagestyle{fancy}
\lhead{\includegraphics[width=7cm]{logoUvA}}
\rhead{\footnotesize \textsc {Report\\ \opdracht}}
\lfoot
{
    \footnotesize \studentA
    \ifthenelse{\isundefined{\studentB}}{}{\\ \studentB}
    \ifthenelse{\isundefined{\studentC}}{}{\\ \studentC}
    \ifthenelse{\isundefined{\studentD}}{}{\\ \studentD}
    \ifthenelse{\isundefined{\studentE}}{}{\\ \studentE}
}
\cfoot{}
\rfoot{\small \textsc {Page \thepage\ of \pageref{LastPage}}}
\renewcommand{\footrulewidth}{0.5pt}

\fancypagestyle{firststyle}
{
    \fancyhf{}
    \renewcommand{\headrulewidth}{0pt}
    \chead{\includegraphics[width=7cm]{logoUvA}}
    \rfoot{\small \textsc {Page \thepage\ of \pageref{LastPage}}}
}

\setlength{\topmargin}{-0.3in}
\setlength{\textheight}{630pt}
\setlength{\headsep}{40pt}

% =================================== DOC INFO ===================================

\newcommand{\titel}{POSIX Threads}
\newcommand{\opdracht}{Assignment 1}
\newcommand{\docent}{Dr. A. Pimentel}
\newcommand{\cursus}{Concurrency and Parallel Programming}
\newcommand{\vakcode}{5062COPP6Y}
\newcommand{\datum}{\today}
\newcommand{\studentA}{Robin Klusman}
\newcommand{\uvanetidA}{10675671}
\newcommand{\studentB}{Maico Timmerman}
\newcommand{\uvanetidB}{10542590}
%\newcommand{\studentC}{Boudewijn Braams}
\newcommand{\uvanetidC}{10401040}
%\newcommand{\studentD}{Govert Verkes}
\newcommand{\uvanetidD}{10211748}
%\newcommand{\studentE}{Naam student 5}
\newcommand{\uvanetidE}{UvAnetID student 5}

% ===================================  ===================================

\begin{document}
\thispagestyle{firststyle}
\begin{center}
    \textsc{\Large \opdracht}\\[0.2cm]
    \rule{\linewidth}{0.5pt} \\[0.4cm]
    {\huge \bfseries \titel}
    \rule{\linewidth}{0.5pt} \\[0.2cm]
    {\large \datum  \\[0.4cm]}

    \begin{minipage}{0.4\textwidth}
        \begin{flushleft}
            \emph{Student:}\\
            {\studentA \\ {\small \uvanetidA \\[0.2cm]}}
            \ifthenelse{\isundefined{\studentB}}{}{\studentB \\ {\small \uvanetidB \\[0.2cm]}}
            \ifthenelse{\isundefined{\studentC}}{}{\studentC \\ {\small \uvanetidC \\[0.2cm]}}
            \ifthenelse{\isundefined{\studentD}}{}{\studentD \\ {\small \uvanetidD \\[0.2cm]}}
            \ifthenelse{\isundefined{\studentE}}{}{\studentE \\ {\small \uvanetidE \\ [0.2cm]}}
        \end{flushleft}
    \end{minipage}
    ~
    \begin{minipage}{0.4\textwidth}
        \begin{flushright}
            \emph{Supervisor:} \\
            \docent \\[0.2cm]
            \emph{Course:} \\
            \cursus \\[0.2cm]
            \emph{Course code:} \\
            \vakcode \\[0.2cm]
        \end{flushright}
    \end{minipage}\\[1 cm]
\end{center}


% =================================== FRONT PAGE ===================================

\tableofcontents
.\\
\begin{center}
    \includegraphics[width=(\textwidth/5*3)]{parallel_tasks}
\end{center}
\clearpage


% =================================== MAIN TEXT ===================================

\section{Introduction}

For this assignment a parallel programming solution needs to be implemented for
two problems, a wave equation simulation and the Sieve of Eratosthenes. For the
wave simulation the user can specify the amount of wave amplitude points, the
amount of steps it needs to simulate and the desired amount of threads. The
program then calculates all the wave values until it has done the specified
amount of steps.\\
For the Sieve of Eratosthenes assignment the program will use the Sieve of
Eratosthenes algorithm to continuously produce prime numbers until a certain
amount of primes have been found. It will do this using multiple threads, thus
utilising parallelism.

\section{Method}
\subsection{Wave Equation Simulation}

First the specified amount of threads need to be created, these threads will
then all start executing their routine. This routine gets a certain range of
amplitude values that it is allowed to calculate, this range is determined by
dividing the amount of amplitude points by the amount of threads, so that each
thread has the same workload. When a thread finishes its calculations it locks
and decrements the variable that is keeping track of the amount of busy threads.
The lock makes sure no two threads try to decrement it at the same time, causing
it to never reach 0.  If this variable reaches 0, the thread that made it 0 will
perform the buffer swap and broadcast a signal to all waiting threads that they
can continue their work. When the specified amount of time steps has been
completed, the threads will all terminate and the final wave is returned.

\subsection{Sieve of Eratosthenes}

To find prime numbers, the Sieve of Eratosthenes algorithm is used. This
multi-threaded program uses this algorithm to find primes by filtering out
non-primes from a constant flow of natural numbers, $n \in \mathbb{N}$, that is
generated by the generator thread. In this implementation we actually opted to
skip every even number, since those are all multiples of 2 and would thus be
filtered out anyway. This will improve the performance of the program a little
because only half as many numbers need to be processed by the first filter and
one less filter is needed. The generated numbers are put into a queue of a fixed
size, from which the next filter will grab its input. If a number reached the
end of the filter chain, that means it has to be a prime, since its not a
multiple of any previously encountered number. This number will then be printed,
and a new filter thread will be created to filter for multiples of this newly
found prime.\\
The queues used to pass numbers between different threads are of a set size, so
that no starvation can occur for threads further down the chain. Reading and
writing to and from the queues is done mutually exclusive, so that no race
conditions can exist. A thread will first lock the input buffer to see if there
is anything in there. If there is, it will grab the first value and unlock the
input. Then the value is processed and if it is not filtered out, the output
buffer is locked to write this value to it. If there is no space in the output
buffer, or if the input buffer is empty, the thread will wait for a signal
before it can continue.

\section{Results}
\subsection{Wave Equation Simulation}
\subsection{Sieve of Eratosthenes}


\section{Discussion}
\subsection{Wave Equation Simulation}
\subsection{Sieve of Eratosthenes}


% =================================== REFERENCES ===================================

%\clearpage
%\bibliographystyle{unsrt}
%\bibliography{bib}

\end{document}
